\section{Executive summary}
\label{abstract}

For CMPU (Customer Must Pickup and Upload) vet diet orders; the
customer is required to pick up the their prescriptions and upload
them to Chewy. This order group suffers from high cancellation rates;
the top cancellation reason begin ``aged order''; an automatic
cancellation that happens if the order has not been shipped or
otherwise been canceled by the customer after ten days from order
creation. Aged order cancellation is followed by explicit
cancellations by the customer (``customer initiated cancellation'' and
``customer request''). Unfortunately, cancellation reasons
do not provide actionable root causes and require further
investigation.

This report focuses on aged order cancellations. Among CMPU vet diet
orders shipped or canceled between May 1, 2026 and July 31, 2026,
38,314 orders were canceled due to aging, with recorded line net sales
of \$3,470,380.19 as an opportunity baseline. These amounts are 9.88\%
of the 387,707 orders and 9.69\% of the \$35,819,402.26 in line net sales
in the no-recorded-cancellation-reason group, respectively (Table
\ref{tab:financial-opportunity-cancellations}). That comparison group
is not a verified shipped-only population, and the aged-order amount
is not an estimate of recoverable revenue. Customer initiated and
customer requested cancellations have a combined recorded line-sales
baseline of \$3,988,013.32.

Within seven days after an aged cancellation, 2,316 of 30,382 customers
(7.62\%) place a qualifying non-cancelled reorder for the same pet and
part number (Table \ref{tab:financial-opportunity-reorder-windows}).
Those reorders have \$221,122.35 in recorded net sales, equivalent to
6.37\% of the \$3,470,380.19 aged-order sales baseline. At 28 days,
recorded follow-up net sales are \$685,508.89, or 19.75\% of that same
baseline. These are observed follow-up sales, not measured incremental
recovery. By 70 days, 7,944 of 30,382 customers (26.15\%) reorder the
same product for the same pet.
Among the 11,550 customers whose first observed customer-upload vet-diet
order aged, 833 (7.21\%) reorder within seven days and 2,509 (21.72\%)
within 70 days. Among the 5,097 customers whose first observed vet-diet
order aged, 800 (15.70\%) reorder within 70 days, using the same pet/product
endpoint (Table \ref{tab:financial-opportunity-new-vet-diet-follow-up}).
For any qualifying Chewy reorder within 70 days, the corresponding
counts are 18,564/30,382 (61.10\%), 6,712/11,550 (58.11\%) and
2,591/5,097 (50.83\%), respectively
(Table \ref{tab:financial-opportunity-reorder-all-chewy}).
The reorder percentage is lowest in the new to vet diet subgroup.


All of the customers seem to have contacted Chewy, most of them at the
creation of the order (Table
\ref{tab:aged-customer-contacts-by-cancellation-day}). Some of them
have contacted Chewy on multiple days (Table
\ref{tab:aged-customer-contact-day-frequency}). People who
contacted closer to the aged order cancellation date have a higher
reorder rate (Table \ref{tab:contact-window-same-pet-part-reorders}).

In the August--October prescription study, 777 of 26,443 aged
first-order lines have a matched upload (2.94\%); 25,666 (97.06\%)
have none in the retained source window (Table
\ref{tab:study-upload-coverage-by-cancellation-reason}).  This
measures matched evidence, not proof that a customer never uploaded.
Compared to the overall upload coverage of 31.18\% , this low coverage
suggests that there is a root cause for aged orders leading to lower
than average upload rates.  Of the 782 distinct matched upload
entities, 52 are currently ACCEPTED (6.65\%); entity acceptance is
distinct from prescription approval.  These snapshot counts allow
uploads and state changes after cancellation.  Of those prescriptions
with uploads, the time between order creation and upload does not
significantly differ from the rest of the study group (Figure
\ref{fig:entity-upload-vs-order-creation}). The accepted aged-order
entities are mostly accepted late; after 10 days (Figure
\ref{fig:entity-upload-state-time-aged}). Approval after acceptance is
not a significant contributor to the age of the order most of the
time; it takes less than a minute after acceptance (95\% of the
time). (Figure
\ref{fig:accepted-upload-to-approval-cumulative}). Finally, Chewy can
expect that if the customer does not upload a prescription in 3 days
after order creation; there is more about 50\% chance that the order
will be canceled as an aged order (Figure
\ref{fig:first-order-line-survival}).

So what can be done to reduce aged orders?  One suggestion to increase
the upload rate (communication with Courtney Gleede) is to send
reminder emails to customer. At this point, Chewy does not contact the
at risk orders' customers. There is no Chewy relevant out-bound
contact data for a comparative/causal analysis. The only contact
happens when the customer is notified that their order has been
canceled.  Assuming that people who receive an email before
cancellation react similarly to the cohort that reorders within 7 days
after cancellation, we expect that out of 30K people considered in the
opportunity study with aged cancellations, 2.3K would react before
order cancellation corresponding to 221K sales. The best effort guess
is that Chewy prevents people from skipping one reorder for this
cohort for an additional gain of 221K net sales.  The upper bound on
the best effort estimate is 390K corresponding to the 14 day reorder
window cohort. For scale, that cohort's \$386,124.68 in observed
follow-up net sales is 1.08\% of the \$35,819,402.26 no-recorded-cancellation-reason
sales baseline in Table \ref{tab:financial-opportunity-cancellations};
this ratio is not a measured incremental sales increase from reminders.

The effect of reminder emails on long term retention is an open
question. Picking up and uploading a prescription is
still an onus with or without a getting a reminder.

